\documentclass{beamer}
\usepackage[utf8]{inputenc}
\usepackage{graphicx}
\usepackage{booktabs}
\usepackage{colortbl}
\usepackage{xcolor}
\usepackage{amsmath}
\usetheme{Madrid}
\usecolortheme{default}

%--------------------
% Title Page Information
%--------------------

\title[ActionCloud]
{ActionCloud: Shared Memory for Teams of AI Agents}

\subtitle{First Review - Phase 1: Building the Plumbing}

\author[Group 14]
{
Group -- 14\\
Naga Shiva -- CB.SC.U4AIE24315\\
Rajashree T -- CB.SC.U4AIE24346\\
Ravishanmugam K -- CB.SC.U4AIE24347\\
Sharvesh -- CB.SC.U4AIE24355
}

\institute[Amrita Vishwa Vidyapeetham]
{
Amrita School of Artificial Intelligence, Amrita Vishwa Vidyapeetham, Coimbatore
}

\date{\today}

\IfFileExists{your-college-logo.png}{\logo{\includegraphics[height=1cm]{your-college-logo.png}}}{}

\begin{document}

% Title Page
\begin{frame}
  \titlepage
\end{frame}

%--------------------
\section{The Problem}
\begin{frame}{The Problem, in Plain Terms}
    \begin{itemize}
        \item Companies now run several AI agents together - one writes code, one does research, one tests, one deploys, one writes docs.
        \item Every time an agent runs, it starts from \textbf{zero}. It has no memory of what it did yesterday, and no idea what another agent already figured out an hour ago.
        \item So the same mistake gets made twice. The same fix gets worked out twice. The same bug gets diagnosed twice.
        \item This gets \textbf{worse}, not better, as more agents and more tasks get added - because nothing connects them.
        \item In one line: \textbf{agents have no shared memory, so they keep redoing each other's work, and that wastes real time and money.}
    \end{itemize}
\end{frame}

%--------------------
\section{Why It Is Hard}
\begin{frame}{Why "Just Add a Shared Database" Is Not Enough}
    \begin{itemize}
        \item The easy fix sounds obvious - let all agents write to one shared memory. But that creates a new danger.
        \item If Agent A solves something \textbf{wrong} and it gets shared blindly, every other agent that reuses it inherits the same mistake. One bad memory can quietly poison the whole team.
        \item So the real question is not just \textbf{"can agents share memory"} - it is \textbf{"what should agents be allowed to trust, and how does something earn that trust"}.
        \item Big cloud providers (AWS, Microsoft) already offer memory storage for agents - but they only solve \textbf{storage}, not \textbf{trust}. They don't decide what should be believed.
        \item \textbf{This project, ActionCloud,} builds that missing trust layer, and proves with a real measured test that it actually saves work and cost - not just claims it.
    \end{itemize}
\end{frame}

%--------------------
\section{The Idea}
\begin{frame}{The Idea, in One Slide}
    \begin{itemize}
        \item Every time an agent finishes a task, it writes down what happened as a small structured note - not a full transcript, just the useful parts: \textbf{what it was asked, what went wrong, what fixed it, and the outcome}.
        \item We call this note an \textbf{Experience}. This is the single most important idea in the whole project - everything else is built around it.
        \item A new experience is not trusted right away. It starts \textbf{private} - visible only to the agent that created it.
        \item It only becomes visible to \textbf{other} agents once it has proven itself - this is what we call \textbf{governance}, explained on the next slides.
        \item Think of it like a new employee's suggestion: it stays a personal note until it's been tried and shown to actually work, before it becomes official team knowledge.
    \end{itemize}
\end{frame}

%--------------------
\section{Related Work}
\begin{frame}{What Already Exists}
\scriptsize
\setlength{\tabcolsep}{3pt}
\renewcommand{\arraystretch}{1.25}

\begin{tabular}{|c|c|p{2.8cm}|p{2.5cm}|p{2.4cm}|p{2.2cm}|}
\hline
\rowcolor{blue!70}
\textcolor{white}{\textbf{\#}} &
\textcolor{white}{\textbf{Year}} &
\textcolor{white}{\textbf{Work}} &
\textcolor{white}{\textbf{What It Does}} &
\textcolor{white}{\textbf{Useful Finding}} &
\textcolor{white}{\textbf{What It Misses}} \\
\hline

\rowcolor{blue!8}
1 & 2026 &
Governed Shared Memory for Multi-Agent Systems (``MemClaw'') &
Shares agent memory using fixed rules about who can see what &
Closest existing system to ours &
Never tested against a no-memory baseline; rules are fixed, not earned \\
\hline

2 & 2026 &
Governed Collaborative Memory as Artificial Selection &
A discussion paper comparing different ways to control shared memory &
Directly asks for the kind of experiment we are running &
Just an argument on paper - no working system built \\
\hline

\rowcolor{blue!8}
3 & 2026 &
Governance Decay under Context Compaction &
Studied 1,323 real agent sessions &
When agents ``compress'' old context, safety rules get silently lost 0\%$\to$30\% of the time &
Shows that today's ``memory'' features are not real governance \\
\hline

4 & 2025 &
LEGOMem (Microsoft Research) &
Lets agents reuse step-by-step procedures from before &
The lead/manager agent benefits most from memory &
Treats every agent role the same; no way to build fleet-wide trust \\
\hline

\end{tabular}
\end{frame}

\begin{frame}{What Already Exists (Continued)}
\scriptsize
\setlength{\tabcolsep}{3pt}
\renewcommand{\arraystretch}{1.25}

\begin{tabular}{|c|c|p{2.8cm}|p{2.5cm}|p{2.4cm}|p{2.2cm}|}
\hline
\rowcolor{blue!70}
\textcolor{white}{\textbf{\#}} &
\textcolor{white}{\textbf{Year}} &
\textcolor{white}{\textbf{Work}} &
\textcolor{white}{\textbf{What It Does}} &
\textcolor{white}{\textbf{Useful Finding}} &
\textcolor{white}{\textbf{What It Misses}} \\
\hline

\rowcolor{blue!8}
5 & 2026 &
Always-On Agents Survey (435 papers) &
Reviews everything published about agents keeping long-term state &
Most work focuses on saving/finding memory - almost nobody studies deciding what to trust or forget &
Just a survey - identifies the gap, doesn't fill it \\
\hline

6 & 2025 &
AWS Bedrock AgentCore Memory &
Amazon's official managed memory service for agents &
Proves storage and sharing memory is now an easy, off-the-shelf cloud feature &
No way to check if a memory is trustworthy before sharing it \\
\hline

\rowcolor{blue!8}
7 & 2025 &
Microsoft Agent Framework &
Bundles memory directly into the agent's runtime &
Convenient built-in state handling &
``Memory'' here is really just squeezing old context - not governed trust \\
\hline

\rowcolor{blue!25}
\multicolumn{2}{|c|}{\textbf{Ours}} &
\textbf{ActionCloud} &
\textbf{Shared memory that must earn trust before being reused} &
\textbf{Measured, not assumed, savings} &
--- \\
\hline

\end{tabular}

\vspace{0.2cm}

\begin{block}{The Gap We Are Filling}
Nobody has combined: (1) a working cloud pipeline for agent memory, (2) trust that is \textbf{earned by results}, not fixed rules, and (3) an actual measured comparison against agents with no memory at all. That combination is what makes this project new.
\end{block}

\end{frame}

%--------------------
\section{What Makes This Different}
\begin{frame}{What Makes This Project Different}
    \textbf{The gap in existing work:}
    \begin{itemize}
        \item Cloud companies solve \textbf{storing} memory, not deciding what's \textbf{worth believing}.
        \item The closest research system uses \textbf{fixed rules} to decide what's shared - and never compares itself to a system with no memory, so nobody knows if it actually helps.
        \item Other work shows that when agents "compress" their context to save space, safety rules quietly disappear - so relying on that as "memory" is risky.
    \end{itemize}

    \vspace{0.25cm}
    \textbf{What we do differently:}
    \begin{itemize}
        \item Trust is \textbf{earned}, not assigned - an experience only becomes shared if it keeps working when other agents try it, not because a rule says so.
        \item We test baseline and memory-enabled agents using the \textbf{exact same code}, with only one setting switched - so if memory helps, we can prove it's really the memory, and not just us writing one version more carefully.
        \item Storing an experience never makes an agent wait - explained on the next slide.
        \item We plan to deliberately plant one bad memory and prove the system catches it - a real test, not just a claim.
    \end{itemize}
\end{frame}

%--------------------
\section{How It Works}
\begin{frame}{How the System Works, Step by Step}
\textbf{Four pieces, and agents only ever talk to the first one:}
\begin{itemize}
    \item \textbf{Agents} - the coding, research, testing, deployment, documentation, and data-analysis agents. They never touch the database directly.
    \item \textbf{One API (front door)} - the only way in or out. This is what makes trust rules actually enforceable - nobody can sneak around it.
    \item \textbf{A queue and background workers} - when an agent saves an experience, it doesn't get processed instantly. It's placed in a line, and a separate background worker picks it up and does the real work.
    \item \textbf{Storage} - PostgreSQL for the structured records, and (from Phase 2) a knowledge graph and a vector search index for smarter retrieval.
\end{itemize}

\vspace{0.2cm}
\textbf{The flow in one line:} \scriptsize Agent $\rightarrow$ API $\rightarrow$ Queue $\rightarrow$ Background Worker $\rightarrow$ Database $\rightarrow$ Search $\rightarrow$ another Agent finds it.

\vspace{0.2cm}\normalsize
\textbf{Why the queue matters:} processing an experience properly (Phase 2 onward) takes a few seconds. If the agent had to wait for that every time, using ActionCloud would make agents \textbf{slower} than having no memory at all - and that would defeat the entire point before we even measured anything. So the agent submits and immediately moves on; the slow work happens quietly in the background.
\end{frame}

%--------------------
\begin{frame}{How Trust Is Earned}
\textbf{An experience is "guilty until proven innocent."} It starts \texttt{PRIVATE}, and can only climb by being reused successfully:
\begin{center}
\scriptsize
\texttt{PRIVATE} $\rightarrow$ \texttt{AGENT-ONLY} $\rightarrow$ \texttt{SHARED} $\rightarrow$ \texttt{PROVEN} $\rightarrow$ \texttt{ORGANIZATION-WIDE}
\end{center}

\begin{itemize}
    \item A component called the \textbf{Memory Judge} decides when to promote an experience - based on who created it and how often it has actually worked when reused. Nothing gets shared just because it exists.
    \item If a memory later turns out to be wrong, it is \textbf{marked as replaced}, never silently deleted - so there's always a paper trail of what was believed and when.
    \item An agent \textbf{cannot} fake its own trust level or confidence score - this isn't a rule we have to remember to enforce, it's built into the data format itself: those fields simply don't exist on what an agent is allowed to submit.
    \item Every promotion or demotion is written to a separate log, so we can always explain \textbf{why} something is trusted.
\end{itemize}
\end{frame}

%--------------------
\section{Our Roadmap}
\begin{frame}{Our Three-Phase Roadmap}
    \begin{itemize}
        \item \textbf{Phase 1 - Build the Plumbing} \hfill \textcolor{blue}{\textbf{[DONE, tested locally]}}
        \begin{itemize}\small
            \item The basic save-and-search loop, a fair way to test with/without memory, 2 of 6 agent types working.
        \end{itemize}
        \item \textbf{Phase 2 - Make It Smart} \hfill \textcolor{gray}{[Next]}
        \begin{itemize}\small
            \item Pull out real knowledge from each experience, build a knowledge graph, add meaning-based search, build the real trust system, add the remaining 4 agent types.
        \end{itemize}
        \item \textbf{Phase 3 - Prove It Works} \hfill \textcolor{gray}{[Final]}
        \begin{itemize}\small
            \item Run 100 real tasks both with and without memory, measure the actual difference, and calculate real cost savings.
        \end{itemize}
    \end{itemize}
    \vspace{0.15cm}
    \footnotesize \textbf{Two rules we locked in from day one:} (1) saving an experience never slows an agent down; (2) the "with memory" and "without memory" agents run on the exact same code, so any difference we measure later is genuinely caused by the memory, not by how carefully we coded each version. \textbf{Long-term goal:} any outside agent framework should be able to plug into ActionCloud through a simple API, not just our own agents.
\end{frame}

%--------------------
\begin{frame}{What We Actually Built in Phase 1}
    \begin{itemize}
        \item \textbf{The API} - the front door. Lets an agent save an experience or search for one. Saving confirms instantly; the real work happens afterward, in the background.
        \item \textbf{The queue} - where a saved experience waits its turn. Built so the exact same code will work on a real cloud queue later, with zero changes.
        \item \textbf{The background worker} - constantly listens to the queue, saves each experience to the database, and decides its starting trust level. Safe to run twice on the same item by accident - it won't create duplicates.
        \item \textbf{The agent library} - a small toolkit so agents don't need to know any of the above; they just call ``search'' and ``save.''
        \item \textbf{The proof script} - an automated check that runs the whole loop end-to-end and confirms it actually works, instead of us just assuming it does.
    \end{itemize}
\end{frame}

%--------------------
\begin{frame}{The Experience Format}
\textbf{Every worker and every future measurement reads from this one shape - it's the foundation everything else stands on:}

\begin{itemize}\small
    \item \textbf{What an agent submits:} what it was asked to do, what problem it hit, what it did, what solved it, the outcome - plus numbers like tokens used, time taken, and cost.
    \item \textbf{What actually gets stored:} the same thing, plus trust information added \textbf{only by the system} - never by the agent - like trust level, confidence, and how many times it's been reused.
    \item \textbf{What gets handed back on search:} a smaller view with just enough to judge whether to trust it - not the full internal history.
\end{itemize}

\vspace{0.2cm}
\textbf{Key point:} an agent physically cannot set its own trust level - that field doesn't even exist on the form it's allowed to fill in. Trust is enforced by the data design itself, not by good behaviour.

\vspace{0.15cm}\footnotesize
We already reserved space in the database for Phase 2 features (search-by-meaning, knowledge graph links) - so adding them later needs no risky changes to data that's already been collected.
\end{frame}

%--------------------
\section{Tech Stack}
\begin{frame}{What We Built It With}
\begin{table}[h]
    \centering
    \small
    \begin{tabular}{ll}
        \toprule
        \textbf{Part} & \textbf{Tool Used} \\
        \midrule
        Front door (API) & FastAPI (Python web framework) \\
        Waiting line (queue) & Amazon SQS (tested locally with a free simulator) \\
        Main storage & PostgreSQL (with vector search built in) \\
        Knowledge graph (Phase 2) & Neo4j, free tier \\
        Talking to the database & Plain SQL - easier to debug than a heavy framework \\
        AI model calls & A free fake model for testing, real model added later \\
        Local development & Docker (runs the whole stack on one laptop) \\
        Testing & pytest (automated checks) \\
        \bottomrule
    \end{tabular}
\end{table}
\vspace{0.15cm}
\footnotesize
\textbf{No AWS account or paid AI key needed yet} - Phase 1 runs entirely for free on a laptop, so we could build and fix the whole loop before spending any real money.
\end{frame}

%--------------------
\section{The Experiment}
\begin{frame}{How We Will Test If This Actually Helps}
\scriptsize
\setlength{\tabcolsep}{4pt}
\renewcommand{\arraystretch}{1.3}
\begin{tabular}{lll}
\toprule
\textbf{} & \textbf{Agents Without Memory} & \textbf{Agents With ActionCloud} \\
\midrule
Starts each task & From zero, every time & By checking for similar past work first \\
Can reuse past fixes & No & Yes \\
Trust checking & --- & Yes, before sharing anything \\
\bottomrule
\end{tabular}

\vspace{0.25cm}\normalsize
\textbf{The plan:} give both groups the exact same 100 tasks (across 5 agent types, 20 each), including tasks that repeat or closely resemble each other on purpose - so we can actually see whether memory catches the repeats. Each task runs once per group, and the memory group only ever learns from what happens \textbf{during this same test} - nothing is fed in beforehand, so the result can't be accidentally rigged.
\end{frame}

%--------------------
\begin{frame}{What We Will Measure}
\scriptsize
\begin{itemize}
    \item \textbf{Success rate} - how often each group actually completes the task correctly.
    \item \textbf{Tokens, tool calls, and time used} - the real cost of each task, compared between groups.
    \item \textbf{Repeated work} - how much of the total work was solving something that was already solved before. Lower is better.
    \item \textbf{Reuse rate} - what share of tasks were actually helped by a memory that already existed.
    \item \textbf{Estimated money saved} - converting the measured token/time savings into an actual dollar figure, using real pricing - not a guess.
\end{itemize}
\vspace{0.2cm}
\normalsize
\textbf{How we'll rank which memory to show first (Phase 2):} a simple scoring formula that rewards a memory for being relevant, having worked before, and being recent - and penalizes it for being expensive to use or getting old and stale.
\[
\text{Score} = \alpha \cdot \text{Relevance} + \beta \cdot \text{Past Success} + \gamma \cdot \text{Freshness} - \delta \cdot \text{Cost} - \varepsilon \cdot \text{Staleness}
\]
\end{frame}

%--------------------
\section{Results So Far}
\begin{frame}{What We Have Proven So Far}
\textbf{All 19 end-to-end checks pass.} The whole stack runs locally using Docker.

\vspace{0.15cm}
\textbf{Confirmed working, not just assumed:}
\begin{itemize}\small
    \item The full loop: an agent saves an experience, and it correctly makes its way through the queue into the database.
    \item The core test that matters most: \textbf{a second, different agent can find and read the first agent's saved experience.}
    \item Accidentally processing the same item twice does not create duplicate data.
    \item All 16 automated unit tests pass.
\end{itemize}

\vspace{0.2cm}
\textbf{Four real bugs we hit and fixed - included honestly, because this is what real engineering looks like:}
\begin{itemize}\footnotesize
    \item A missing setting almost made the system quietly try to connect to \textbf{real AWS} instead of our free local test setup - fixed by loading the settings file properly and blocking that possibility outright.
    \item Our database container was accidentally colliding with a Postgres already installed on a laptop - fixed by moving it to a different port.
    \item \textbf{A logical trap:} new experiences start private, and private experiences are correctly hidden from other agents - but with nothing to promote them, nothing could \textbf{ever} become shared. Fixed by adding a basic promotion rule.
    \item A bug where a shared tool object was mixing up \textit{which agent} an experience actually belonged to - fixed by always passing identity explicitly.
\end{itemize}
\end{frame}

%--------------------
\begin{frame}{Progress: Done, In Progress, and Next}
\small
\begin{columns}[T]

\begin{column}{0.32\textwidth}
\textbf{Phase 1} \\
\footnotesize The Plumbing \\
\textcolor{blue}{\textbf{DONE (local)}}
\rule{\linewidth}{0.4pt}
\begin{itemize}\footnotesize
\item[\checkmark] Experience format + database
\item[\checkmark] Save/search API
\item[\checkmark] Queue + background worker
\item[\checkmark] Agent toolkit
\item[\checkmark] Fair with/without-memory test setup
\item[\checkmark] 2 of 6 agent types + full proof test
\end{itemize}
\end{column}

\begin{column}{0.32\textwidth}
\textbf{Phase 2} \\
\footnotesize Making It Smart \\
\textcolor{gray}{\textbf{NEXT}}
\rule{\linewidth}{0.4pt}
\begin{itemize}\footnotesize
\item[$\circ$] Pull real knowledge out automatically
\item[$\circ$] Build the knowledge graph
\item[$\circ$] Meaning-based search
\item[$\circ$] Real trust system
\item[$\circ$] Remaining 4 agent types
\item[$\circ$] Open connector for outside agents
\end{itemize}
\end{column}

\begin{column}{0.32\textwidth}
\textbf{Phase 3} \\
\footnotesize Proving It Works \\
\textcolor{gray}{\textbf{FINAL}}
\rule{\linewidth}{0.4pt}
\begin{itemize}\footnotesize
\item[$\circ$] 100-task real test
\item[$\circ$] With vs.\ without memory run
\item[$\circ$] Real measurements + cost savings
\item[$\circ$] Dashboard + final report
\end{itemize}
\end{column}

\end{columns}

\vspace{0.2cm}\footnotesize
\textbf{Still left in Phase 1:} move from our local test setup to the real cloud (AWS) - the code itself won't change, only which server it's pointed at. A spending alert will be set up first, before anything real is created.
\end{frame}

%--------------------
\begin{frame}{How We Compare to What Already Exists}

\begin{itemize}
    \item Comparing ActionCloud against the closest existing systems:
\end{itemize}

\vspace{0.3cm}

\begin{table}[h!]
\centering
\renewcommand{\arraystretch}{1.4}
\resizebox{\textwidth}{!}{%
\begin{tabular}{llllll}
\toprule
\textbf{System} & \textbf{Cloud Pipeline} & \textbf{Trust Checking} & \textbf{Trust Is Earned} & \textbf{Tested vs.\ No-Memory} & \textbf{Real Cost Numbers} \\
\midrule
AWS AgentCore Memory        & Yes & No       & No  & No  & No \\
Microsoft Agent Framework   & Yes & Partial\footnotemark[1] & No  & No  & No \\
MemClaw (closest research)  & Yes & Yes      & No (fixed rules) & No  & No \\
LEGOMem (Microsoft)         & Partial & No   & No  & No  & No \\
\midrule
\textbf{ActionCloud} & \textbf{Yes} & \textbf{Yes} & \textbf{Yes} & \textbf{Planned (Phase 3)} & \textbf{Planned (Phase 3)} \\
\bottomrule
\end{tabular}%
}
\end{table}
\footnotetext[1]{Only through compressing old context - which has been shown to quietly lose safety rules, not enforce them.}

\end{frame}

%--------------------
\section{Conclusion}
\begin{frame}{Where We Stand and What's Next}
    \textbf{Where we stand:}
    \begin{itemize}
        \item \textbf{Phase 1 is complete and verified} - the basic save-and-find loop genuinely works, proven by an automated test, not just by us saying so.
        \item The design choice to make saving \textbf{never slow an agent down} works exactly as planned.
        \item Trust rules are built into the system itself, not just something we promise to follow.
        \item Because the with-memory and without-memory agents share the exact same code, any result we measure later will be trustworthy.
    \end{itemize}

    \vspace{0.25cm}
    \textbf{What's next:}
    \begin{itemize}
        \item \textbf{Phase 2:} teach the system to pull out real knowledge automatically, build the knowledge graph, add meaning-based search, and build the real trust system.
        \item \textbf{Phase 3:} run the actual 100-task experiment and get real, honest numbers - not projections.
        \item \textbf{Move to real AWS} - mostly a configuration change, not a rewrite.
        \item Replace our current placeholder trust rule and our current "every task counts as a success" stand-in before any Phase 3 result can be trusted.
    \end{itemize}
\end{frame}

%--------------------
\section{References}
\begin{frame}{References}

\begin{thebibliography}{99}
\footnotesize

\bibitem{1} AWS Documentation, ``Get started with AgentCore Memory,'' Amazon Bedrock AgentCore Developer Guide, 2025.

\bibitem{2} Amazon Web Services, Inc., ``Amazon Bedrock AgentCore - FAQs,'' aws.amazon.com/bedrock/agentcore/faqs/, 2025.

\bibitem{3} Y. Margalit, N. Cohen-Inger, E. Avram, R. Taig, O. Margalit, ``Governed Shared Memory for Multi-Agent LLM Systems,'' arXiv:2606.24535, 2026.

\bibitem{4} Cuadros, Maiga, Meskhidze, Curtis-Trudel, ``Governed Collaborative Memory as Artificial Selection in LLM-Based Multi-Agent Systems,'' arXiv:2605.04264, 2026.

\bibitem{5} S. Chen, ``Governance Decay,'' arXiv:2606.22528, 2026.

\bibitem{6} T. Ding, A. Nannapaneni, B. Liu, L. Zhang, ``Always-On Agents: A Survey,'' arXiv:2606.30306, 2026.

\bibitem{7} D. Han, C. Couturier, D. Madrigal Diaz, X. Zhang, V. R\"uhle, S. Rajmohan, ``LEGOMem: Procedural Memory for Multi-Agent LLM Systems,'' Microsoft Research, arXiv:2510.04851, 2025.

\bibitem{8} ``Graph-based Agent Memory: Taxonomy, Techniques, and Applications,'' arXiv preprint, 2026.

\bibitem{9} Ramakrishnan, et al., ``pgvector: Open-source vector similarity search for Postgres,'' 2024.

\bibitem{10} Amazon Web Services, ``Amazon Simple Queue Service Developer Guide,'' 2024.

\end{thebibliography}

\end{frame}

\end{document}
